\documentclass[12pt]{article}
\usepackage[utf8]{inputenc}
\usepackage[T1]{fontenc}
\usepackage{hyperref}
\title{Gbaya-Mbodomo Language Resources for K-12 Curriculum Development}
\author{Compiled by Claude Sonnet 5.5 (claude-sonnet-5-5) for Andr\'as Kornai}
\date{}
\begin{document}
\maketitle

\section{Introduction}
Gbaya-Mbodomo (ISO 639-3 gmm; Glottocode gbay1281) is a Gbaya language of Cameroon, with about 8,000 speakers according to a 1992 figure cited by Wikipedia; other sources give larger numbers. It was split from Southwest Gbaya [gso] because intelligibility was too low. Materials found in search are very limited.

\section{Primary Educational Resources}
No primers, textbooks or school programmes in Gbaya-Mbodomo were found in search results.

\section{Linguistic and Reference Materials}
\begin{itemize}
\item Overview: \url{https://en.wikipedia.org/wiki/Mbodomo_language}.
\item Glottolog record of the split from Southwest Gbaya: \url{https://glottolog.org/resource/languoid/id/sout3222}.
\item Huma-Num CRDO language record; it lists a grammar by Virginia Boyd (1997): \url{https://cocoon.huma-num.fr/exist/crdo/lexvo.xql?code=gmm}.
\item SIL Cameroon archive entries (reported to include a 1996 French-language lexicon of the Gandima Tongo dialect and a folk tale text; contents not individually verified): \url{https://www.silcam.org/resources/archives/32272}, \url{https://www.silcam.org/resources/archives/47623}, \url{https://www.silcam.org/resources/archives/47475}.
\end{itemize}

\section{Religious and Community Resources}
Joshua Project profile of the Gbaya Mbadomo people: \url{https://joshuaproject.net/people_groups/19813}.

\section{Digital and Audio Resources}
Global Recordings Network lists no recordings in Gbaya-Mbodomo itself: \url{https://www.globalrecordings.net/en/language/gmm}. Recordings exist in the related Baya: Doka variety: \url{https://globalrecordings.net/en/language/1594}.

\section{Curriculum Development Context}
Little exists. A mother-tongue curriculum would have to start from new orthography and materials development with the community, using the SIL lexicon and Boyd's grammar as linguistic bases.
\end{document}